%% Author Contributions Checklist, generated by make_jasa.py from ACC_form.md.
\documentclass[11pt]{article}
\usepackage[margin=1in]{geometry}
\usepackage[T1]{fontenc}
\usepackage{amssymb}
\usepackage{xurl}
\usepackage[colorlinks=true,urlcolor=blue,linkcolor=black]{hyperref}
\setlength{\emergencystretch}{3em}
\providecommand{\tightlist}{\setlength{\itemsep}{0pt}\setlength{\parskip}{0pt}}
\setcounter{secnumdepth}{0}
\title{Author Contributions Checklist (ACC) Form\\[6pt]
\large Betting against the simulator: Monte Carlo e-processes under adaptive actions}
\date{}
\begin{document}
\maketitle
This form documents the artifacts associated with the article (i.e., the
data and code supporting the computational findings) and describes how
to reproduce the findings.

\hypertarget{part-1-data}{%
\section{Part 1: Data}\label{part-1-data}}

\begin{itemize}
\tightlist
\item[$\boxtimes$]
  This paper does not involve analysis of external data (i.e., no data
  are used or the only data are generated by the authors via simulation
  in their code).
\item[$\boxtimes$]
  I certify that the author(s) of the manuscript have legitimate access
  to and permission to use the data used in this manuscript.
\end{itemize}

\hypertarget{abstract}{%
\subsection{Abstract}\label{abstract}}

All data are generated by the code. The simulation studies and the
multi-arm multi-stage trial simulate their data; the trial uses the
patient parameters of simglucose, the open-source implementation of the
UVA/Padova simulator. The language-model audit uses replies generated by
the code from the publicly released Qwen2.5-1.5B-Instruct (8, 4 and 3
bits) and Qwen2.5-0.5B-Instruct (4 bits) models: 1500 to 3000 replies
per model and prompt, for four prompts, stored in
\texttt{data/llm\_audit/replies.json}.

\hypertarget{availability}{%
\subsection{Availability}\label{availability}}

\begin{itemize}
\tightlist
\item[$\boxtimes$]
  Data \textbf{are} publicly available.
\item[$\square$]
  Data \textbf{cannot be made} publicly available.
\end{itemize}

If the data are publicly available:

\begin{itemize}
\tightlist
\item[$\boxtimes$]
  Data are available online at:
  \url{https://github.com/elenabortolato/evalues4generative}, file
  \texttt{data/llm\_audit/replies.json} (TO CONFIRM: for double-blind
  review, an anonymous mirror of the repository can replace this link).
\item[$\square$]
  Data are available as part of the paper's supplementary material.
\item[$\square$]
  Data are publicly available by request.
\item[$\square$]
  Data are available through some other mechanism.
\end{itemize}

\hypertarget{file-formats}{%
\subsection{File format(s)}\label{file-formats}}

\begin{itemize}
\tightlist
\item[$\boxtimes$]
  CSV or other plain text: JSON.
\item[$\square$]
  Software-specific binary format.
\item[$\square$]
  Standardized binary format.
\item[$\square$]
  Other.
\end{itemize}

\hypertarget{data-dictionary}{%
\subsection{Data dictionary}\label{data-dictionary}}

\begin{itemize}
\tightlist
\item[$\boxtimes$]
  Provided by authors in the following file(s): \texttt{data/README.md},
  section ``Language-model replies''.
\item[$\square$]
  Data file(s) is (are) self-describing.
\item[$\square$]
  Available at the following URL.
\end{itemize}

\hypertarget{additional-information-optional}{%
\subsection{Additional Information
(optional)}\label{additional-information-optional}}

The replies are generated by
\texttt{experiments/exp16\_llm\_generate.py} with one seed per reply;
the model files are the public GGUF releases named in
\texttt{data/README.md}.

\hypertarget{part-2-code}{%
\section{Part 2: Code}\label{part-2-code}}

\hypertarget{abstract-1}{%
\subsection{Abstract}\label{abstract-1}}

The Python package \texttt{e4g} implements the Monte Carlo e-process,
the hedged wealth and the score learners: the Bayesian learning rule on
the logistic loss and on the simulation likelihood, predictive-first
updates and a Bayesian-bootstrap forest. The folder
\texttt{experiments/} contains one script per experiment, which writes
its numbers to \texttt{results/} (JSON), its figures to
\texttt{figures/} and, for the language-model audit, the LaTeX of Tables
6 and 7. The concurrent methods compared in the paper (Forré and
Brenner; SLICE) are reimplemented in the same scripts. The folder
\texttt{tests/} checks, among other things, that every bet has mean
exactly one under the null.

\hypertarget{code-formats}{%
\subsection{Code format(s)}\label{code-formats}}

\begin{itemize}
\tightlist
\item[$\boxtimes$]
  Script files: Python.
\item[$\boxtimes$]
  Package: Python (\texttt{e4g}, used from the repository root, not
  installed).
\item[$\square$]
  Reproducible report.
\item[$\square$]
  Shell script.
\item[$\square$]
  Other.
\end{itemize}

\hypertarget{supporting-software-requirements}{%
\subsection{Supporting software
requirements}\label{supporting-software-requirements}}

\hypertarget{primary-software-used}{%
\subsubsection{Primary software used}\label{primary-software-used}}

\begin{itemize}
\tightlist
\item
  Python 3.11 for all analyses.
\item
  Python 3.13 with llama-cpp-python 0.3.36, only to regenerate the
  language-model replies, which are stored in the repository.
\end{itemize}

\hypertarget{libraries-and-dependencies-used-by-the-code}{%
\subsubsection{Libraries and dependencies used by the
code}\label{libraries-and-dependencies-used-by-the-code}}

numpy 2.4.6, scipy 1.17.1, scikit-learn 1.9.1, pandas 3.0.6, matplotlib
3.11.2, simglucose 0.2.11 (installed with \texttt{-\/-no-deps}; only its
patient parameters are used), pytest. Playwright with Chromium is used
only to draw the scheme of Figure 1. The list is in
\texttt{REPRODUCIBILITY.md} (Section 1) and \texttt{requirements.txt}.

\hypertarget{supporting-systemhardware-requirements-optional}{%
\subsection{Supporting system/hardware requirements
(optional)}\label{supporting-systemhardware-requirements-optional}}

A Linux machine with 4 CPU cores; no GPU. Regenerating the replies needs
the GGUF model files (1 to 2 GB each).

\hypertarget{parallelization-used}{%
\subsection{Parallelization used}\label{parallelization-used}}

\begin{itemize}
\tightlist
\item[$\square$]
  No parallel code used.
\item[$\boxtimes$]
  Multi-core parallelization on a single machine/node: 4 cores. The seed
  of each replication is fixed by its index, so the results do not
  depend on the number of workers.
\item[$\square$]
  Multi-machine/multi-node parallelization.
\end{itemize}

\hypertarget{license}{%
\subsection{License}\label{license}}

\begin{itemize}
\tightlist
\item[$\square$]
  MIT License (default) (TO CONFIRM: the repository has no LICENSE file
  yet; add one before submission).
\item[$\square$]
  BSD
\item[$\square$]
  GNU GPL v3.0
\item[$\square$]
  Creative Commons
\item[$\square$]
  Other.
\end{itemize}

\hypertarget{additional-information-optional-1}{%
\subsection{Additional information
(optional)}\label{additional-information-optional-1}}

Version of the code used for the submitted manuscript: commit
\texttt{3b9a61cc7373} (TO CONFIRM: replace with the commit or tag of the
submitted version).

\hypertarget{part-3-reproducibility-workflow}{%
\section{Part 3: Reproducibility
workflow}\label{part-3-reproducibility-workflow}}

\hypertarget{scope}{%
\subsection{Scope}\label{scope}}

The provided workflow reproduces:

\begin{itemize}
\tightlist
\item[$\boxtimes$]
  Any numbers provided in text in the paper.
\item[$\boxtimes$]
  The computational method(s) presented in the paper (i.e., code is
  provided that implements the method(s)).
\item[$\boxtimes$]
  All tables and figures in the paper.
\item[$\square$]
  Selected tables and figures in the paper.
\end{itemize}

\hypertarget{location}{%
\subsection{Location}\label{location}}

\begin{itemize}
\tightlist
\item[$\square$]
  As part of the paper's supplementary material.
\item[$\boxtimes$]
  In this Git repository:
  \url{https://github.com/elenabortolato/evalues4generative} (commit as
  above).
\item[$\square$]
  Other.
\end{itemize}

\hypertarget{formats}{%
\subsection{Format(s)}\label{formats}}

\begin{itemize}
\tightlist
\item[$\square$]
  Single master code file.
\item[$\square$]
  Wrapper (shell) script(s).
\item[$\square$]
  Self-contained literate programming file.
\item[$\boxtimes$]
  Text file (e.g., a readme-style file) that documents workflow:
  \texttt{REPRODUCIBILITY.md}.
\item[$\square$]
  Makefile.
\item[$\square$]
  Other.
\end{itemize}

\hypertarget{instructions}{%
\subsection{Instructions}\label{instructions}}

Install the packages listed in Section 1 of \texttt{REPRODUCIBILITY.md}
and run \texttt{python\ -m\ pytest\ tests}. \texttt{REPRODUCIBILITY.md}
then lists, for every table and figure of the paper and of the
supplement, the command that produces it, the output file and the
running time; all commands are run from the repository root. The JSON
files in \texttt{results/} hold the numbers reported in the paper, so
each table can be checked against them without re-running, and
\texttt{experiments/exp16\_tables.py} prints Tables 6 and 7 from them.
Every replication is seeded. On a different processor, OpenBLAS can
round matrix inversions differently, which can move a few entries of
Tables 6 and 7 within Monte Carlo error; \texttt{REPRODUCIBILITY.md}
documents this.

\hypertarget{expected-run-time}{%
\subsection{Expected run-time}\label{expected-run-time}}

Approximate time needed to reproduce the analyses on a standard desktop
machine:

\begin{itemize}
\tightlist
\item[$\square$]
  \textless{} 1 minute
\item[$\square$]
  1-10 minutes
\item[$\square$]
  10-60 minutes
\item[$\boxtimes$]
  1-8 hours: about 5 hours on 4 cores for the main paper, using the
  stored language-model replies. Regenerating the replies takes about 3
  more hours, and the supplement about 8 more hours.
\item[$\square$]
  \textgreater{} 8 hours
\item[$\square$]
  Not feasible to run on a desktop machine.
\end{itemize}

\hypertarget{additional-information-optional-2}{%
\subsection{Additional information
(optional)}\label{additional-information-optional-2}}

\texttt{examples/llm\_audit\_minimal.py} runs one step of the
language-model audit and prints each quantity of the bet, in a few
seconds.

\hypertarget{notes-optional}{%
\section{Notes (optional)}\label{notes-optional}}

The paper is reviewed double-blind. The links above identify the author;
for the review copy they can be replaced by an anonymous mirror of the
repository (for example through Anonymous GitHub), to be replaced by the
public repository on acceptance.

\end{document}
